\section{System Overview}
\label{sec:system}

\textsc{taxonomy\_agent} is structured as an orchestrator LLM that holds
the running taxonomy as state and a parallel judge LLM that scores
sampled items against that taxonomy each iteration. The orchestrator
interacts with the corpus exclusively through six tools; all taxonomy
mutations go through a typed revise DSL that is validated before being
applied. The loop terminates when the judge's don't-fit rate falls
below a threshold, subject to a minimum iteration floor.
\todo{Insert system diagram (figures/architecture.pdf).}

\subsection{Tools}
\label{sec:system:tools}

The orchestrator's action space is the following six tools:
\begin{itemize}
  \item \texttt{sample\_items} -- draw a batch of unlabelled documents.
  \item \texttt{get\_taxonomy} -- read the current taxonomy state.
  \item \texttt{revise\_taxonomy} -- apply a typed list of revise ops.
  \item \texttt{classify\_with\_judge} -- dispatch the current batch to
        the parallel judge for scoring.
  \item \texttt{propose\_novelties\_with\_judge} -- ask the judge to
        surface candidate categories from don't-fit items.
  \item \texttt{finalize\_classify} -- run the judge over the full
        corpus once convergence is declared.
\end{itemize}
\todo{Flesh out each tool's I/O signature and prompt template.}

\subsection{Revise DSL}
\label{sec:system:revise}

All taxonomy mutations are expressed as a list of typed operations:
\textsc{add}, \textsc{rename}, \textsc{edit}, \textsc{merge},
\textsc{split}, \textsc{drop}. Each op is validated against the current
taxonomy (target labels must exist, new labels must not collide,
splits must specify children) before any mutation is applied; an
invalid batch is rejected as a whole so the taxonomy is never left in
a partially-edited state.
\todo{Table of op signatures and validation rules; worked example of a
two-op merge+rename batch.}

\subsection{Convergence}
\label{sec:system:convergence}

Convergence is decided by the judge's don't-fit rate over the last
sampled batch, gated by a configurable minimum-iteration floor to
prevent premature termination on small early samples. An escape hatch
allows the orchestrator to terminate explicitly when it judges further
edits unproductive.
\todo{State the threshold formula, default floor, and behaviour when
the orchestrator triggers the escape hatch.}
